\documentclass{article}

% packages
\usepackage{amsmath, amsthm, thmtools, amsfonts, amssymb, luacode, catchfile, tikzducks, hyperref, ifthen}
\ifcsname c@kobocompile\endcsname
	\usepackage[a5paper, total={1072pt, 1448pt}, margin=10pt, includeheadfoot]{geometry} % set page margins
\else
	\usepackage[a4paper, margin=50pt, includeheadfoot]{geometry}
\fi
\usepackage[shortlabels]{enumitem}
\usepackage[skip=3pt, indent=0pt]{parskip}

% language
\usepackage[bidi=basic, layout=tabular, provide=*]{babel}
\ifcsname c@english\endcsname
	\babelprovide[main, import]{english}
\else
	\babelprovide[main, import]{hebrew}
	\babelprovide{rl}
\fi
%\babelfont{rm}{Libertinus Serif}
\babelfont{rm}[Renderer=Harfbuzz]{Libertinus Serif}
\babelfont{sf}{Libertinus Sans}
\babelfont{tt}{Libertinus Mono}

% style
\AddToHook{cmd/section/before}{\clearpage}	% Add line break before section
\linespread{1.3}
\setcounter{secnumdepth}{0}		% Remove default number tags from sections, this won't do well with theorems
\AtBeginDocument{\setlength{\belowdisplayskip}{3pt}}
\AtBeginDocument{\setlength{\abovedisplayskip}{3pt}}
\graphicspath{ {../images/} }

% operators
\DeclareMathOperator\cis{cis}
\DeclareMathOperator\Sp{Sp}
\DeclareMathOperator\tr{tr}
\DeclareMathOperator\im{Im}
\DeclareMathOperator\re{Re}
\DeclareMathOperator\diag{diag}
\DeclareMathOperator*\lowlim{\underline{lim}}
\DeclareMathOperator*\uplim{\overline{lim}}
\DeclareMathOperator\rng{rng}
\DeclareMathOperator\Sym{Sym}
\DeclareMathOperator\Arg{Arg}
\DeclareMathOperator\Log{Log}
\DeclareMathOperator\dom{dom}
\DeclareMathOperator\supp{Supp}
\DeclareMathOperator\var{Var}
\DeclareMathOperator\cov{Cov}

% commands
%\renewcommand\qedsymbol{\textbf{מש''ל}}
%\renewcommand\qedsymbol{\fbox{\emoji{lizard}}}
\newcommand{\Aa}[0]{\mathcal{A}}
\newcommand{\Bb}[0]{\mathcal{B}}
\newcommand{\CC}[0]{\mathbb{C}}
\newcommand{\Cc}[0]{\mathcal{C}}
\newcommand{\EE}[0]{\mathbb{E}}
\newcommand{\FF}[0]{\mathbb{F}}
\newcommand{\Ff}[0]{\mathcal{F}}
\newcommand{\Ii}[0]{\mathcal{I}}
\newcommand{\Gg}[0]{\mathcal{G}}
\newcommand{\Ll}[0]{\mathcal{L}}
\newcommand{\Mm}[0]{\mathcal{M}}
\newcommand{\NN}[0]{\mathbb{N}}
\newcommand{\Nn}[0]{\mathcal{N}}
\newcommand{\PP}[0]{\mathbb{P}}
\newcommand{\Pp}[0]{\mathcal{P}}
\newcommand{\QQ}[0]{\mathbb{Q}}
\newcommand{\RR}[0]{\mathbb{R}}
\newcommand{\Rr}[0]{\mathcal{R}}
\newcommand{\Ss}[0]{\mathcal{S}}
\newcommand{\TT}[0]{\mathbb{T}}
\newcommand{\Uu}[0]{\mathcal{U}}
\newcommand{\Vv}[0]{\mathcal{V}}
\newcommand{\Ww}[0]{\mathcal{W}}
\newcommand{\ZZ}[0]{\mathbb{Z}}
\newcommand{\acts}[0]{\circlearrowright}
\newcommand{\explain}[2] {
	\begin{flalign*}
		 && \text{#2} && \text{#1}
	\end{flalign*}
}
\newcommand{\maketitleprint}[0]{ \begin{center}
	%\begin{tikzpicture}[scale=3]
	%	\duck[graduate=gray!20!black, tassel=red!70!black]
	%\end{tikzpicture}	
	\includegraphics[width=6cm]{cover}
\end{center}
}

% theorem commands
\newtheoremstyle{c_remark}
	{}	% Space above
	{}	% Space below
	{}% Body font
	{}	% Indent amount
	{\bfseries}	% Theorem head font
	{}	% Punctuation after theorem head
	{.5em}	% Space after theorem head
	{\thmname{#1}\thmnumber{ #2}\thmnote{ \normalfont{\text{(#3)}}}}	% head content
\newtheoremstyle{c_definition}
	{3pt}	% Space above
	{3pt}	% Space below
	{}% Body font
	{}	% Indent amount
	{\bfseries}	% Theorem head font
	{}	% Punctuation after theorem head
	{.5em}	% Space after theorem head
	{\thmname{#1}\thmnumber{ #2}\thmnote{ \normalfont{\text{(#3)}}}}	% head content
\newtheoremstyle{c_plain}
	{3pt}	% Space above
	{3pt}	% Space below
	{\itshape}% Body font
	{}	% Indent amount
	{\bfseries}	% Theorem head font
	{}	% Punctuation after theorem head
	{.5em}	% Space after theorem head
	{\thmname{#1}\thmnumber{ #2}\thmnote{ \text{(#3)}}}	% head content

\ifcsname c@english\endcsname
	\theoremstyle{plain}
	\newtheorem{theorem}{Theorem}[section]
	\newtheorem{lemma}[theorem]{Lemma}
	\newtheorem{proposition}[theorem]{Proposition}
	\newtheorem*{proposition*}{Proposition}
	%\newtheorem{corollary}[theorem]{אין חלופה עברית}

	\theoremstyle{definition}
	\newtheorem{definition}[theorem]{Definition}
	\newtheorem*{definition*}{Definition}
	\newtheorem{example}{Example}[section]
	\newtheorem{exercise}{Exercise}[section]

	\theoremstyle{remark}
	\newtheorem*{remark}{Remark}
	\newtheorem*{solution}{Solution}
	\newtheorem{conclusion}[theorem]{Conclusion}
	\newtheorem{notation}[theorem]{Notation}
\else
	\theoremstyle{c_plain}
	\newtheorem{theorem}{משפט}[section]
	\newtheorem{lemma}[theorem]{למה}
	\newtheorem{proposition}[theorem]{טענה}
	\newtheorem*{proposition*}{טענה}
	%\newtheorem{corollary}[theorem]{אין חלופה עברית}

	\theoremstyle{c_definition}
	\newtheorem{definition}[theorem]{הגדרה}
	\newtheorem*{definition*}{הגדרה}
	\newtheorem{example}{דוגמה}[section]
	\newtheorem{exercise}{תרגיל}[section]

	\theoremstyle{c_remark}
	\newtheorem*{remark}{הערה}
	\newtheorem*{solution}{פתרון}
	\newtheorem{conclusion}[theorem]{מסקנה}
	\newtheorem{notation}[theorem]{סימון}
\fi

% Questions related commands
\newcounter{question}
\setcounter{question}{1}
\newcounter{sub_question}
\setcounter{sub_question}{1}

\ifcsname c@english\endcsname
	\newcommand{\question}[1][0]{
		\ifthenelse{#1 = 0}{}{\setcounter{question}{#1}}
		\section{Question \arabic{question}}
		\addtocounter{question}{1}
		\setcounter{sub_question}{1}
	}

	\newcommand{\subquestion}[1][0]{
		\ifthenelse{#1 = 0}{}{\setcounter{sub_question}{#1}}
		\subsection{Part \alph{sub_question}}
		\addtocounter{sub_question}{1}
	}
\else
	\newcommand{\question}[1][0]{
		\ifthenelse{#1 = 0}{}{\setcounter{question}{#1}}
		\section{שאלה \arabic{question}}
		\addtocounter{question}{1}
		\setcounter{sub_question}{1}
	}

	\newcommand{\subquestion}[1][0]{
		\ifthenelse{#1 = 0}{}{\setcounter{sub_question}{#1}}
		\subsection{סעיף \localecounter{letters.gershayim}{sub_question}}
		\addtocounter{sub_question}{1}
	}
\fi

% import lua and start of document
\directlua{common = require ('../common')}

\GetEnv{AUTHOR}

% headers
\author{\AUTHOR}
\date\today

\title{פתרון מטלה 10 --- מבוא לטופולוגיה, 80516}
% chktex-file 9
% chktex-file 17

\usepackage{pgfplots}
\pgfplotsset{width=7cm,compat=1.18}

\begin{document}
\maketitle
\maketitleprint[purple]

\question{}
נניח ש־$0 < n \in \NN$,
נראה כי התנאים הבאים שקולים,
\begin{enumerate}
	\item לכל $f : S^n \to \RR^n$ רציפה קיים $x \in S^n$ כך ש־$f(x) = f(-x)$
	\item לא קיימת פונקציה רציפה אי־זוגית $S^n \to S^{n - 1}$
\end{enumerate}
\begin{proof}
	$1 \implies 2$,
	נניח בשלילה שקיימת $g : S^n \to S^{n - 1}$ אי־זוגית, קרי $f(-x) = -f(x)$ לכל $x \in S^n$.
	תהי פונקציה המהווה הרחבה רציפה של $g$, $f : S^n \to \RR^n$ (קיימת מהלמה של אוריסון), אז קיימת $x \in S^n$ כך ש־$f(x) = f(-x)$, אבל $f \restriction S^n = g$ ולכן $g(x) \in S^{n - 1}$ ובפרט $g(x) = g(-x) = -g(x)$.
	נסיק אם כך ש־$g(-x) = g(x) = 0$ בלבד, אבל $0 \notin S^{n - 1}$ בסתירה לקיום $g$.

	$2 \implies 1$,
	תהי $f : S^n \to \RR^n$.
	נגדיר,
	\[
		g(x)
		= \frac{f(x) - f(-x)}{2}
	\]
	נבחין כי הפונקציה $g$ היא אי־זוגית.
	אילו $g(x) = 0$ ל־$x$ כלשהו, אז $f(x) = f(-x)$ וסיימנו, לכן נניח אחרת.
	לכן בפרט נוכל להגדיר,
	\[
		h(x)
		= \frac{g(x)}{\lVert g(x) \rVert}
	\]
	ונקבל פונקציה $h : S^n \to S^{n - 1}$ רציפה, ובהכרח גם אי־זוגית, בסתירה ישירה להנחה שלנו, ולכן $h$ לא קיימת, ובהתאם $g$ בעלת שורש.
\end{proof}

\question{}
נניח את משפט בורסוק־אולם ונוכיח שאין שיכון $S^n \hookrightarrow \RR^n$.
\begin{proof}
	נניח בשלילה ש־$f : S^n \hookrightarrow \RR^n$ שיכון.
	אז קיים $x \in S^n$ כך ש־$f(x) = f(-x)$, אבל מחד־חד ערכיות השיכון $x \ne -x \implies f(x) \ne f(-x)$ ולכן $x = 0$ בלבד, אבל $0 \notin S^n$, וזו סתירה להנחה שלנו.
\end{proof}

\question{}
נגדיר את $S^1 = \{ z \in \CC \mid |z| = 1 \}$ ויהי פולינום $p(z) = z^n + a_{n - 1} z^{n - 1} + \cdots + a_0$.

\subquestion{}
אם $q : \CC \to \CC$ רציפה כך ש־$|z| = r \implies q(z) \ne 0$ עבור $r \ge 0$ כלשהו.
אז הפונקציה,
\[
	f_r^q : I \to S^1,
	\quad
	f_r^q(s)
	= \frac{q(r e^{2 \pi i s}) / q(r)}{|q(r e^{2 \pi i s}) / q(r)|}
\]
היא פונקציה רציפה.

\subquestion{}
יהי $R > |a_{n - 1}| + \cdots + |a_0| + 1$,
לכל $0 \le t \le 1$ נגדיר פולינום,
\[
	p_t(z)
	= z^n + t(a_{n - 1} z^{n - 1} + \cdots + a_0)
\]
ונגדיר את הפונקציה $h : I^2 \to S^1$ המוגדרת על־ידי $h(t, s) = f_R^{p_t}(s)$.
נראה כי $h$ היא הומוטופיה מ־$s \mapsto e^{2 \pi i n s}$ ל־$f_R^p$.
\begin{proof}
	ברור כי $h$ רציפה, ולכן עלינו רק לבדוק את התנאים להומוטופיה.
	\[
		h(0, s)
		= f_R^{z^n}(s)
		= \frac{R^n e^{2 \pi i n s} / R^n}{|R^n e^{2 \pi i n s} / R^n|}
		= \frac{e^{2 \pi i n s}}{|e^{2 \pi i n s}|}
		= e^{2 \pi i n s}
	\]
	מהצד השני,
	\[
		h(1, s)
		= f_R^p
	\]
	ישירות מאיך שהגדרנו את $p_t$ וההתלכדות $p_1 = p$.
	לבסוף גם,
	\[
		h(t, 0)
		= f_R^{p_t}(0)
		= \frac{p_t(R) / p_t(R)}{|p_t(R) / p(R)|}
		= 1
	\]
	ובאופן דומה נקבל גם $h(t, 1) = 1$.
\end{proof}

\subquestion{}
נסיק שאם $n \ge 1$ אז $f_R^p$ לא הומוטופית ללולאה קבועה.
\begin{proof}
	הומוטופיה היא יחס טרנזיטיבי, לכן מספיק שנראה שאין הומוטופיה בין הלולאה $t \mapsto e^{2 \pi i n t}$ ללולאה הקבועה $0$.
	למעשה זו טענה שהוכחה בכיתה על־ידי שימוש בכיסוי של המעגל על־ידי הישר הממשי.
\end{proof}

\subquestion{}
נסיק את המשפט היסודי של האלגברה.
\begin{proof}
	נניח בשלילה שהוא לא מתקיים, כלומר קיים פולינום ממעלה חיובית אשר אין לו שורש.
	נסמן פולינום זה כ־$p$, ונבחין כי $f_R^p$ מוגדרת לכל $s$, ונגדיר,
	\[
		h(t, s)
		= f_{\gamma(s)}^p(t)
	\]
	עבור מסילה $\gamma = [0, R]$.
	זוהי כמובן פונקציה רציפה ישירות מהגדרה ולכן יש הומוטופיה בין $f_0^p$ לבין $f_R^p$.
	אבל זו סתירה לסעיף הקודם, ולכן נסיק שלא קיים פולינום כזה.
\end{proof}

\end{document}
